\documentclass[11pt,a4paper]{article}
\usepackage{amsmath,amssymb,amsthm}
\usepackage{enumitem}
\usepackage[margin=2.4cm]{geometry}
\usepackage[colorlinks=true,linkcolor=blue,citecolor=blue,urlcolor=blue]{hyperref}

\theoremstyle{plain}
\newtheorem{lemma}{Lemma}[section]
\newtheorem{proposition}[lemma]{Proposition}
\newtheorem{theorem}[lemma]{Theorem}

\newcommand{\Stab}{\operatorname{Stab}}
\newcommand{\im}{\operatorname{im}}
\newcommand{\cof}{\operatorname{cof}}
\newcommand{\id}{\mathrm{id}}
\newcommand{\Set}{\mathbf{Set}}

\title{Stabilization Points of $\omega$-Accessible Set Functors: $\Stab(F) \in \{0, 1, \omega, \omega + \omega\}$}
\date{}

\begin{document}
\maketitle

\section*{0. Preliminaries}

Let $F : \Set \to \Set$ be an endofunctor. Its terminal sequence $(A_\alpha, f_\alpha^\beta : A_\beta \to A_\alpha)_{\alpha \le \beta}$ is defined by transfinite recursion:
\begin{itemize}[leftmargin=2em]
  \item $A_0 = 1$ (singleton set);
  \item $A_{\alpha+1} = F(A_\alpha)$ and $f_{\alpha+1}^{\beta+1} = F(f_\alpha^\beta)$;
  \item for limit ordinals $\lambda$, $A_\lambda = \varprojlim_{\alpha < \lambda} A_\alpha$ is the limit of compatible families, with projections $f_\alpha^\lambda$.
\end{itemize}
We say $F$ \textbf{stabilizes at $\gamma$} if $p_\gamma := f_\gamma^{\gamma+1} : A_{\gamma+1} \to A_\gamma$ is bijective (which then persists to all $\delta \ge \gamma$); the least stabilizing ordinal is $\Stab(F)$.

$F$ is \textbf{$\omega$-accessible} (finitary) if for every set $X$ and every $u \in F(X)$, there is a finite subset $S \subseteq X$ with $u \in \im(F(i_S))$, where $i_S : S \hookrightarrow X$ is the inclusion. For $z_0 \in X$, we also write $i_{\{z_0\}} : 1 \to X$ for the constant map $* \mapsto z_0$.

\section{Four Basic Lemmas}

\begin{lemma}[Trnková]
Let $F : \Set \to \Set$ and $A, B \subseteq X$. If $A \cap B \neq \emptyset$, then
\[
  \im(F(i_A)) \cap \im(F(i_B)) = \im(F(i_{A \cap B})).
\]
\end{lemma}

\begin{proof}
($\supseteq$) Immediate from $i_{A \cap B} = i_A \circ i_{A\cap B \hookrightarrow A} = i_B \circ i_{A\cap B \hookrightarrow B}$.

($\subseteq$) Pick $c \in A \cap B$ and define retractions $e_A, e_B : X \to X$ by
\[
  e_A(x) = \begin{cases} x, & x \in A, \\ c, & x \notin A, \end{cases}
  \qquad
  e_B(x) = \begin{cases} x, & x \in B, \\ c, & x \notin B. \end{cases}
\]
Then $e_A \circ i_A = i_A$, $e_B \circ i_B = i_B$, and $e_A \circ e_B = i_{A \cap B} \circ q$ for the map $q : X \to A \cap B$ sending $X \setminus (A \cap B)$ to $c$. If $t \in \im(F(i_A)) \cap \im(F(i_B))$, writing $t = F(i_A)(u_A) = F(i_B)(u_B)$ gives $F(e_A)(t) = t$ and $F(e_B)(t) = t$. Hence
\[
  t = F(e_A)\bigl(F(e_B)(t)\bigr) = F(e_A \circ e_B)(t) = F(i_{A \cap B})\bigl(F(q)(t)\bigr) \in \im(F(i_{A \cap B})). \qedhere
\]
\end{proof}

\begin{lemma}
Let $F : \Set \to \Set$ and $A, B \subseteq X$ be disjoint ($A \cap B = \emptyset$). If $x \in \im(F(i_A)) \cap \im(F(i_B))$, then for every $z_0 \in X$, $x \in \im(F(i_{\{z_0\}}))$.
\end{lemma}

\begin{proof}
Fix $z_0 \in X$ and define $f : X \to X$ by $f(y) = z_0$ for $y \in A$ and $f(y) = y$ for $y \notin A$.
\begin{enumerate}[label=(\arabic*)]
  \item Since $A \cap B = \emptyset$, $f \circ i_B = i_B$, so $x \in \im(F(i_B))$ implies $F(f)(x) = x$.
  \item Meanwhile $f \circ i_A = i_{\{z_0\}} \circ {!}_A$ where ${!}_A : A \to 1$. Writing $x = F(i_A)(a)$ gives
    \[
      x = F(f)(x) = F(f \circ i_A)(a) = F(i_{\{z_0\}})\bigl(F({!}_A)(a)\bigr) \in \im(F(i_{\{z_0\}})). \qedhere
    \]
\end{enumerate}
\end{proof}

\begin{lemma}
Let $F : \Set \to \Set$. Suppose there exist a set $X$ and maps $q : X \to A_1$, $s : A_1 \to X$ with $q \circ s = \id_{A_1}$, and suppose $F({!}_X) : F(X) \to F(1)$ is bijective. Then $p_1 : F(A_1) \to A_1$ is bijective, so $\Stab(F) \le 1$.
\end{lemma}

\begin{proof}
From $q \circ s = \id_{A_1}$ we get $F(q) \circ F(s) = \id_{F(A_1)}$, so $F(s)$ is injective. Also ${!}_X \circ s = {!}_1$ and ${!}_1 \circ q = {!}_X$.
\begin{itemize}[leftmargin=2em]
  \item Since $F(!_X) = F(!_1) \circ F(q)$ is surjective, $F(!_1)$ is surjective;
  \item Since $F(!_1) = F(!_X) \circ F(s)$ is a composite of injections, $F(!_1)$ is injective.
\end{itemize}
Thus $F(!_1) : F(A_1) \to F(1)$ is bijective, which identifies with $p_1 : A_2 \to A_1$ via $A_1 \cong F(1)$.
\end{proof}

\begin{lemma}
If $A_1 = F(1) \neq \emptyset$, then for every ordinal $\alpha \ge 1$ there exists a section $s_\alpha : A_1 \to A_\alpha$ such that $f_1^\alpha \circ s_\alpha = \id_{A_1}$. In particular, $|A_\alpha| \ge |A_1|$ and every $A_\alpha \neq \emptyset$.
\end{lemma}

\begin{proof}
Each $c \in F(1) = A_1$ defines a constant coalgebra $1 \xrightarrow{c} F(1)$, which generates by transfinite recursion a unique compatible family of elements $s_\alpha(c) \in A_\alpha$ pushed forward via constant maps $1 \to A_\delta$ at successor stages and projection limits at limit stages, satisfying $f_1^\alpha(s_\alpha(c)) = c$ (Worrell 2005). Hence $s_\alpha$ is injective and $|A_\alpha| \ge |A_1|$.
\end{proof}

\section{Excluding Non-Trivial Successor Convergence Points}

\begin{proposition}
Let $F : \Set \to \Set$ be $\omega$-accessible and $|X| \ge 2$. If there exists a proper subset $T \subsetneq X$ such that $F(i_T) : F(T) \to F(X)$ is surjective, then for every $z \in X$, both $F(i_{\{z\}}) : F(1) \xrightarrow{\sim} F(X)$ and $F({!}_X) : F(X) \xrightarrow{\sim} F(1)$ are bijective.
\end{proposition}

\begin{proof}
\textbf{Step 1}: Choose $g \in X \setminus T$. Since $T \subseteq X \setminus \{g\}$, surjectivity of $F(i_T)$ makes $F(i_{X\setminus\{g\}})$ surjective. Conjugating by the swap bijection $(g\; b) : X \xrightarrow{\sim} X$ shows $F(i_{X\setminus\{b\}})$ is surjective for every $b \in X$.

\textbf{Step 2}: Fix $z \in X$ and $x \in F(X)$; we show $x \in \im(F(i_{\{z\}}))$:
\begin{itemize}[leftmargin=2em]
  \item \textbf{Case 1 ($|X| = M$ finite, $M \ge 2$)}: Since $\{z\} = \bigcap_{b \in X \setminus \{z\}} (X \setminus \{b\})$ with every finite sub-intersection containing $z$, repeated Lemma~1.1 yields $\im(F(i_{\{z\}})) = \bigcap_{b \neq z} \im(F(i_{X\setminus\{b\}})) = F(X)$.
  \item \textbf{Case 2 ($|X|$ infinite)}: Pick a finite support $B_0 \subseteq X$ for $x$ by $\omega$-accessibility. If $B_0 = \emptyset$, then $\emptyset \hookrightarrow X$ factors through $i_{\{z\}}$, so $x \in \im(F(i_{\{z\}}))$. If $B_0 \neq \emptyset$, then $X \setminus B_0 = \bigcap_{b \in B_0} (X \setminus \{b\})$ is infinite (hence nonempty); iterated Lemma~1.1 over the finite set $B_0$ gives $x \in \im(F(i_{X \setminus B_0}))$, and Lemma~1.2 yields $x \in \im(F(i_{\{z\}}))$.
\end{itemize}

\textbf{Step 3}: From $1 \xrightarrow{i_{\{z\}}} X \xrightarrow{{!}_X} 1$ we have $F({!}_X) \circ F(i_{\{z\}}) = \id_{F(1)}$. Together with Step~2 (surjectivity of $F(i_{\{z\}})$), $F(i_{\{z\}})$ is bijective, and therefore so is its inverse $F({!}_X)$.
\end{proof}

\begin{theorem}
Let $F : \Set \to \Set$ be $\omega$-accessible. Then $\Stab(F) \neq \alpha + 1$ for every ordinal $\alpha \ge 1$.
\end{theorem}

\begin{proof}
Suppose toward contradiction that $\Stab(F) = \alpha + 1$ with $\alpha \ge 1$.
\begin{itemize}[leftmargin=2em]
  \item If $A_1 = \emptyset$, existence of $A_2 \to A_1 = \emptyset$ forces $A_2 = \emptyset$, so $p_1 : \emptyset \to \emptyset$ is bijective while $p_0 : \emptyset \to 1$ is not, yielding $\Stab(F) = 1$, impossible.
  \item If $|A_1| = 1$, $p_0 : A_1 \to A_0$ is already bijective, yielding $\Stab(F) = 0$, impossible.
\end{itemize}
Hence $|A_1| \ge 2$, and by Lemma~1.4, $|A_\delta| \ge |A_1| \ge 2$ for every $\delta \ge 1$.

Because $\Stab(F) = \alpha + 1$, $p_{\alpha+1} = F(p_\alpha)$ is bijective while $p_\alpha : A_{\alpha+1} \to A_\alpha$ is not bijective.

\textbf{Case 1 ($p_\alpha$ not surjective)}: Set $T := \im(p_\alpha) \subsetneq A_\alpha$. Factoring $p_\alpha = i_T \circ q$, surjectivity of $F(p_\alpha)$ gives surjectivity of $F(i_T)$. Since $|A_\alpha| \ge 2$, Proposition~2.1 on $X = A_\alpha$ gives $F(A_\alpha) \cong F(1)$, and Lemmas~1.3 and~1.4 yield $\Stab(F) \le 1$, contradicting $\alpha + 1 > 1$.

\textbf{Case 2 ($p_\alpha$ surjective, not injective)}: Choose $y_0 \in A_\alpha$ and $x_1 \neq x_2 \in A_{\alpha+1}$ with $p_\alpha(x_1) = p_\alpha(x_2) = y_0$. Take a right inverse $\sigma_0$ of $p_\alpha$ and define sections $\sigma_1, \sigma_2 : A_\alpha \to A_{\alpha+1}$ differing only at $y_0$, with $\sigma_k(y_0) = x_k$. Since $p_\alpha \circ \sigma_k = \id_{A_\alpha}$ and $F(p_\alpha) = p_{\alpha+1}$ is bijective, $F(\sigma_1) = F(\sigma_2) = p_{\alpha+1}^{-1}$ are bijective. Let $T_0 := \sigma_1(A_\alpha \setminus \{y_0\})$ and $S_k := \im(\sigma_k) = T_0 \cup \{x_k\}$. Because $|A_\alpha| \ge 2$, $T_0 \neq \emptyset$; since $x_1 \notin T_0$, $T_0 \subsetneq A_{\alpha+1}$ and $S_1 \cap S_2 = T_0$. Surjectivity of $F(\sigma_k)$ implies surjectivity of $F(i_{S_k})$, and Lemma~1.1 yields surjectivity of $F(i_{T_0})$. Applying Proposition~2.1 to $X = A_{\alpha+1}$ and $T_0$, then Lemmas~1.3 and~1.4, again yields $\Stab(F) \le 1$, contradiction.
\end{proof}

\section{Worrell's Upper Bound and the 4-Point Classification}

\begin{theorem}[Worrell's Upper Bound]
Every $\omega$-accessible set functor $F : \Set \to \Set$ stabilizes by $\omega + \omega$.
\end{theorem}

\begin{proof}
If $A_1 = \emptyset$, $F$ stabilizes at $1 \le \omega + \omega$. Otherwise all $A_\alpha$ are nonempty by Lemma~1.4, and $\omega$-accessibility implies tail monicity: $f_\alpha^\beta : A_\beta \to A_\alpha$ is injective for all $\omega \le \alpha \le \beta$. In particular, $p_{\omega+\omega}$ is injective.

To show $p_{\omega+\omega}$ is surjective, let $y \in A_{\omega+\omega}$ and put $y_1 := f_{\omega+1}^{\omega+\omega}(y) \in A_{\omega+1} = F(A_\omega)$. For every $k \in \mathbb{N}$, naturality gives $y_1 = F(f_\omega^{\omega+k})\bigl(f_{\omega+k+1}^{\omega+\omega}(y)\bigr) \in \im(F(i_{C_k}))$, where $C_k := \im(f_\omega^{\omega+k}) \subseteq A_\omega$ forms a descending chain. By injectivity of each $f_\omega^{\omega+k}$, any element of $\bigcap_{k \in \mathbb{N}} C_k$ lifts coherently to $A_{\omega+\omega}$, so $\bigcap_{k \in \mathbb{N}} C_k = C_\omega := \im(f_\omega^{\omega+\omega})$.

By $\omega$-accessibility, pick a finite support $S \subseteq A_\omega$ for $y_1$. Choosing $N \in \mathbb{N}$ above the first exit index of every $s \in S \setminus C_\omega$ ensures $S \cap C_N \subseteq C_\omega$. Either $S \cap C_N \neq \emptyset$ (Lemma~1.1) or $S \cap C_N = \emptyset$ (Lemma~1.2 with $z_0 = f_\omega^{\omega+\omega}(y) \in C_\omega$); in both cases $y_1 \in \im(F(i_{C_\omega}))$. Inverting the bijection $A_{\omega+\omega} \xrightarrow{\sim} C_\omega$ gives $x \in F(A_{\omega+\omega}) = A_{\omega+\omega+1}$ with $f_{\omega+1}^{\omega+\omega+1}(x) = y_1 = f_{\omega+1}^{\omega+\omega}(y)$. Finally, injectivity of $f_{\omega+1}^{\omega+k+1}$ forces $f_{\omega+k+1}^{\omega+\omega+1}(x) = f_{\omega+k+1}^{\omega+\omega}(y)$ for every $k \in \mathbb{N}$, hence $x$ and $y$ agree on all coordinates $< \omega + \omega$, giving $p_{\omega+\omega}(x) = y$.
\end{proof}

\begin{theorem}[Four-Point Convergence Classification]
Every $\omega$-accessible set functor $F : \Set \to \Set$ has a unique exact stabilization ordinal $\Stab(F)$, and
\[
  \Stab(F) \in \{0,\; 1,\; \omega,\; \omega + \omega\}.
\]
\end{theorem}

\begin{proof}
By Theorem~3.1, $\gamma := \Stab(F) \le \omega + \omega$. Theorem~2.2 rules out all finite ordinals $n \ge 2$ and all ordinals $\omega + (n + 1)$ ($n \in \mathbb{N}$) strictly between $\omega$ and $\omega + \omega$, leaving precisely $\{0, 1, \omega, \omega + \omega\}$.
\end{proof}

\end{document}
